% ECONOMICS_PROBLEM: {"id": "Q54-176", "title": "Economic consequences of climate tipping points", "jel_code": "Q54", "source_id": "OP-0431"}
% !TeX program = xelatex
\documentclass[11pt,a4paper]{article}
\usepackage[margin=23mm]{geometry}
\usepackage{fontspec}
\setmainfont{Latin Modern Roman}
\usepackage{amsmath,amssymb,mathtools}
\usepackage{xcolor,enumitem,booktabs,longtable,array,xurl,hyperref}
\definecolor{muted}{HTML}{53636C}
\hypersetup{colorlinks=true,urlcolor=blue,linkcolor=blue}
\setlength{\parindent}{0pt}
\setlength{\parskip}{5pt}
\newcommand{\R}{\mathbb R}
\newcommand{\E}{\mathbb E}
\newcommand{\N}{\mathbb N}
\newcommand{\ind}{\mathbf 1}
\newcommand{\Var}{\operatorname{Var}}
\newcommand{\Cov}{\operatorname{Cov}}
\newcommand{\supp}{\operatorname{supp}}
\DeclareMathOperator*{\argmin}{arg\,min}
\DeclareMathOperator*{\argmax}{arg\,max}
\newcommand{\entrymeta}[3]{{\sffamily\small\color{muted}\textbf{\texttt{#1}}\quad|\quad #2\par #3\par}}
\newcommand{\Problemref}[1]{\texttt{#1}}
\newcommand{\JELref}[2]{\texttt{#2} (JEL \texttt{#1})}
\begin{document}
\section*{Economic consequences of climate tipping points}
\textbf{Problem ID:} \texttt{Q54-176}\quad \textbf{JEL:} \texttt{Q54}\par
% BEGIN ECONOMICS STATEMENT
\label{problem:OP-0431}
{\sffamily\small\color{muted}Original subject group: Climate, environment and energy\par}
\label{definitions:problem:30007713}
\textbf{Audit classification:} Background; proposed extension. Specified irreversible hazard timing, population weights, equilibrium selection and model-set bounds.
\subsection*{Definitions and precise research target}
\textit{Editorial formalization.} Consider a calibrated world economy with a finite set of regions \(r\), annual periods \(t=0,\ldots,100\), and one specified physical tipping event, such as a major ice-sheet transition. Let \(T_t\) be global mean temperature, \(S_t\in\{0,1\}\) the irreversible event indicator, and \(q(T_t)\) its conditional transition probability when \(S_t=0\). Consumption \(C_{rt}>0\), adaptation investment \(A_{rt}\geq0\), capital, and migration are equilibrium outcomes. A stated damage map links the event to sea-level rise and regional production. For discount factor \(\beta\in(0,1)\), positive regional welfare weights \(w_r\), and increasing concave utility \(u_r\), define the incremental expected welfare loss
\[L=E\!\left[\sum_{t=0}^{100}\beta^t\sum_r w_r\{u_r(C_{rt}^{0})-u_r(C_{rt}^{1})\}\right].\]
Superscripts denote otherwise identical economies respectively excluding and including the event risk; the expectation integrates physical and economic shocks. Determine an identified or transparently bounded range for \(L\) using defensible transition probabilities and endogenous adaptation. Separate uncertainty about physical hazards from uncertainty about economic propagation, and test whether allowing migration, insurance, and correlated regional damages reverses policy rankings. This is a proposed quantitative model comparison for one event, rather than a claim that tipping points lack existing economic analyses.

Use a filtered probability space and a physical state \(X_t\). Define \(P(S_{t+1}=1\mid\mathcal F_t,S_t=0)=q(X_t)\in[0,1]\) and \(S_{t+1}=1\) whenever \(S_t=1\); temperature-only hazards are a restriction. Declare carbon dynamics, sea-level mapping, production, household budgets, adaptation technology, and a measurable equilibrium-selection rule. Couple the risk/no-risk experiments by common exogenous shocks, setting the hazard to zero in regime 0. Regional consumption is per capita, with fixed person weights or explicitly modeled population weights; utility sums must use the same population accounting convention. Require \(E\sum_t\beta^t\sum_r w_r|u_r(C_{rt}^a)|<\infty\). For admissible calibrated models \(\mathcal M\), the target interval is \([\inf_{M\in\mathcal M}L(M),\sup_{M\in\mathcal M}L(M)]\); neither sign nor narrowness is assumed. All expectations use a stated baseline population and probability law. Identification means every admissible data-generating model with the same observed-data law yields the same displayed target; otherwise report the identified set. Exact equations, horizons and assumptions are editorial, rather than source-given canonical conjectures.
\subsection*{Variants and assumption changes}
\begin{enumerate}
\item \textbf{Editorial interacting-events variant.} Replace \(S\) by \(S\in\{0,1\}^J\) and specify joint transition probabilities, allowing one event to change another's hazard. Compare total \(L\) with summed single-event losses; equality is not assumed.
\item \textbf{Editorial ambiguity variant.} Replace \(P\) by a declared nonempty set \(\mathcal P\) and report \(\inf_{P\in\mathcal P}L_P,\sup_{P\in\mathcal P}L_P\), separating hazard and equilibrium-selection ambiguity.
\end{enumerate}
\subsection*{References and source audit}
\textbf{1.} Climate macroeconomics survey; no exact fifty-year or tipping-event conjecture verified. Locator: NBER working-paper record and abstract; AEA forthcoming record, DOI 10.1257/jel.20261831. Primary indexed record verified; direct NBER/PDF retrieval failed. AEA forthcoming abstract inspected; no final journal publication year asserted. URL: \url{https://www.nber.org/papers/w33567}.

\textbf{2.} SCC synthesis discusses damage persistence, tipping, and nonmarket substitution; proposed estimands are editorial. Locator: Introduction; Section 1.1; Section 2. Full text or primary publication page inspected. URL: \url{https://pmc.ncbi.nlm.nih.gov/articles/PMC11670083/}.
\begin{enumerate}\raggedright
\item\hypertarget{definitions:source:30007713:1}{} Adrien Bilal; James H. Stock. 2025. \emph{A Guide to Macroeconomics and Climate Change}. NBER working-paper record and abstract; AEA forthcoming record, DOI 10.1257/jel.20261831.
\url{https://www.nber.org/papers/w33567}.
DOI: \url{https://doi.org/10.3386/w33567}.
\item\hypertarget{definitions:source:30007713:2}{} Frances C. Moore; Moritz A. Drupp; James Rising; Simon Dietz; Ivan Rudik; Gernot Wagner. 2024. \emph{Synthesis of evidence yields high social cost of carbon due to structural model variation and uncertainties}. Introduction; Section 1.1; Section 2.
\url{https://pmc.ncbi.nlm.nih.gov/articles/PMC11670083/}.
DOI: \url{https://doi.org/10.1073/pnas.2410733121}.
\end{enumerate}

\subsection*{Common conventions: Source mathematical conventions}

These conventions apply to each entry unless explicitly replaced.

\textbf{Models and identification.} A primitive model \(m\) specifies all probability laws, feasible choices, information, equilibrium and measurement rules used by its target. The admissible class \(\mathcal M\) is fixed independently of outcomes. If \(P_O(m)\) is its observable probability law and \(\tau(m)\) the target, its identified set at observable law \(P\) is
\[
 \mathcal I_\tau(P)=\{\tau(m):m\in\mathcal M,\ P_O(m)=P\}.
\]
Point identification means that this set is a singleton. An empty set signals incompatibility of data and assumptions. Coordinate bounds need not describe the joint set. Bounds are sharp only if every retained target is attainable or arbitrarily approachable by compatible models. Finite bounds require explicit restrictions on unobserved quantities; unrestricted confounding, hidden wealth, extrapolation or arbitrary equilibrium selection can yield uninformative sets.

\textbf{Causal interventions.} A potential outcome \(Y_i(p)\) is the outcome under a complete policy regime \(p\), including timing, financing, information and market spillovers. The target population and units are fixed before treatment. With interference, use the complete assignment vector or a justified exposure mapping; individual no-interference notation alone cannot identify an economy-wide intervention. A difference of mean potential outcomes does not require the cross-world joint distribution. Conditioning on post-treatment migration, survival, enrollment or employment changes the estimand and needs separate justification.

\textbf{Statistical statements.} An estimator, confidence set, forecast or test specifies a sampling law, model class and loss/coverage criterion. Uniform \((1-\alpha)\)-coverage means \(\inf_{m\in\mathcal M}\Pr_m\{\tau(m)\in C_n\}\geq1-\alpha\), or an explicitly asymptotic version. The whole parameter space trivially has coverage, so informativeness requires a stated width, risk or power objective. A forecasting improvement is relative to a named benchmark, information set and evaluation population.

\textbf{Equilibrium and optimization.} An equilibrium comprises feasible plans, optimality, beliefs where needed, budget constraints and all market/resource clearing. The primitives or a stated benchmark define each correspondence \(\mathcal E(p;m)\). Nonemptiness is assumed or proved, never inferred just from compact policy sets. A selected-equilibrium target uses a declared selection rule; a robust target quantifies over all permitted equilibria. Use suprema or infima unless attainment is established. An \(\varepsilon\)-optimal policy is feasible and within \(\varepsilon\) of the finite optimum value. Report an empty feasible set separately.

\textbf{Welfare and accounting.} Normative utility, interpersonal normalization, social weights, numeraire, discounting and the use of public receipts are primitives. Behavioral utility need not coincide with the welfare criterion. Household welfare includes ownership income and rents once; adding profits again to compensating variation generally double-counts them. A monetary equivalent is defined by an explicit transfer and price convention, with monotonicity and range conditions. Average effects, mechanism contributions, transfers and welfare losses are different targets. Infinite sums require convergence and dynamic feasible plans require terminal or no-Ponzi conditions.

\textbf{Source versions.} A working-paper date, revision date and journal publication year are distinguished. A DOI for a journal version is not automatically the DOI of an earlier manuscript. Locators refer to the stated version and printed pages where specified. A metadata-only check does not verify a claim about a full-text conclusion. Every entry's source subsection narrows the provenance claim accordingly.
% END ECONOMICS STATEMENT
\end{document}
